% =============================================================================
\subsection{Metodologia Experimental e Protocolo}
\label{sec:metodologia}

\subsubsection{Bases de dados}

As três bases (\cref{tab:bases}) são fornecidas já divididas em treino e
teste, com a saída $y$ na última coluna. O treino é usado para evoluir as
soluções e o teste é usado \emph{apenas} ao final, para avaliar o melhor
indivíduo de cada execução. As bases sintéticas possuem duas variáveis de
entrada ($x_0, x_1$); a base \code{concrete1} representa um problema de
dados empíricos multivariados com oito atributos contínuos ($x_0, \dots, x_7$),
totalizando 1030 observações particionadas em 824 amostras de treino e 206
de teste, cuja variável alvo $y$ apresenta magnitude média substancialmente
mais elevada ($\bar{y} \approx 35.76$ e $\sigma_y \approx 16.27$).

\begin{table}[htbp]
  \centering
  \caption{Bases de dados utilizadas. $\sigma_y$ é o desvio-padrão da saída no
           conjunto de teste, usado como referência de escala do RMSE.}
  \label{tab:bases}
  \small
  \begin{tabular}{@{}lcccccl@{}}
    \toprule
    Base & Entradas & Treino & Teste & $\bar{y}_{\text{teste}}$ & $\sigma_{y,\text{teste}}$ & Tipo \\
    \midrule
    \code{synth1}    & 2 & 60  & 600  & 17.75 & 25.89 & Sintética (polinomial) \\
    \code{synth2}    & 2 & 300 & 1000 & 0.36  & 4.03  & Sintética (racional) \\
    \code{concrete1} & 8 & 824 & 206  & 35.76 & 16.27 & Empírica / Real \\
    \bottomrule
  \end{tabular}
\end{table}

Para facilitar a comparação entre bases, também se reporta o coeficiente de
determinação no teste, $R^2 = 1-\mathrm{RMSE}^2/\sigma_{y}^2$.

\subsubsection{Protocolo Experimental e Métricas}

\begin{itemize}
  \item \textbf{Repetições:} cada configuração foi executada 30 vezes com as
        sementes $0,\dots,29$. Como as mesmas sementes são usadas em todas as
        configurações, as amostras podem ser pareadas por semente.
  \item \textbf{Medidas:} RMSE de treino e de teste do melhor indivíduo
        (média $\pm$ desvio-padrão, mediana, intervalo interquartil (IQR),
        mínimo e máximo), tamanho do melhor indivíduo (nós) e tempo de
        execução. Como as distribuições de RMSE são fortemente assimétricas,
        com caudas longas, a \textbf{mediana} é a medida central privilegiada
        na discussão.
  \item \textbf{Testes estatísticos:} teste U de
        Mann--Whitney~\cite{mann1947test} (bilateral, $\alpha=0.05$) para
        comparar configurações; na Fase~2 também o teste de postos sinalizados
        de Wilcoxon~\cite{wilcoxon1945individual}, pareado por semente, e a
        contagem de vitórias/derrotas (W/L) por semente.
  \item \textbf{Um fator por vez:} ao variar um parâmetro, todos os demais
        são mantidos constantes.
\end{itemize}

\subsubsection{Organização das Duas Fases Experimentais}

\begin{description}
  \item[Fase~1] (\cref{sec:fase1}). Configuração base, conforme o guia do
        enunciado: \code{pop}=100, \code{gen}=100, profundidade máxima 7,
        \emph{ramped half-and-half}, torneio com $k=2$, $\pc=0.9$,
        $\pmut=0.05$, mutação de subárvore, elitismo com 1 indivíduo e
        $\mathcal{F}=\{+,-,\times,\div\}$. O estudo é feito no \code{synth1},
        e a configuração resultante é replicada no \code{synth2} e no
        \code{concrete1}. Ao todo, foram 22 configurações $\times$ 30
        execuções.
  \item[Fase~2] (\cref{sec:fase2}). Parte de um \emph{estado atual}
        (\cfg{S0}), igual à configuração base, mas com $\eps$-lexicase (o
        vencedor da Fase~1), avaliado nas três bases. Seis hipóteses (H1--H6)
        são testadas \textbf{isoladamente} contra \cfg{S0} e, em seguida,
        \textbf{empilhadas} em duas ``escadas'' cumulativas (A e B). Todas
        as 14 configurações foram executadas nas três bases
        ($14\times3\times30 = 1260$ execuções), o que permite avaliar também
        o efeito de cada intervenção nas bases para as quais ela
        \emph{não} foi concebida.
\end{description}
